\section{Introduction}
Marking an exam booklet requires more than assigning a total. A teacher must associate each response with a question, recognise partial understanding, justify deductions and decide when an uncertain answer needs another look. Handwriting and multilingual text add recognition and interpretation errors before a marking decision is made. A fluent explanation is insufficient if it cannot be traced to the student's work or the teacher's rubric.

Auto-Assessment treats the rubric as an explicit contract. It accepts an exam specification containing questions, criterion descriptions and model answers, together with typed sheets, scans or photographs. It returns criterion marks, verbatim evidence, actionable feedback, an HTML report and a red-ink annotated booklet. The intended role is selective auto-marking with teacher review, preserving an inspectable basis for both credit and referral.

This report evaluates that role on typed SPRAG responses and a small real handwritten classroom set. The central analytical question is whether review mechanisms improve agreement or merely change its direction. An advocate can recover missed credit; a strict critic can only remove it. The ablation finds that the critic increases under-marking, motivating an advocate-only default chosen on development data. A separate confidence analysis asks how much marking can be accepted automatically at a stated error level.

The contributions are an integrated full-sheet workflow with evidence checks and bounded awards; an empirical comparison of review directions; a confidence-based referral mechanism; and a report that distinguishes typed grading, handwritten end-to-end performance and targeted injection checks. Feedback remains the educational motivation, but its quality has not yet been rated by students and no learning effect is claimed. The measurements support assisted marking within the tested settings, rather than a validated replacement for teachers.
